\hypertarget{appendix-17917d9fdbe3}{}
\section*{Ripple-carry addition}
\chapter{Ripple-carry addition}\label{appendix-17917d9fdbe3-ripple-carry-addition}

\textbf{Curriculum depth:} Introduced · \textbf{First seen:} C002 · \textbf{Latest development:} C002

\section{Definition and mechanism}\label{appendix-17917d9fdbe3-definition-and-mechanism}

Each full adder feeds its carry to the next position, making worst-case circuit propagation depth grow with word width.

Connect carry-out at bit i to carry-in at i+1. For \texttt{01111111\ +\ 1}, a carry traverses the low seven stages. Circuit delay grows with width in this simple design; fixed-width software addition still counts as O(1) with respect to input size.

\section{Prerequisites}\label{appendix-17917d9fdbe3-prerequisites}

\begin{itemize}
\tightlist
\item
  \hyperref[appendix-d1a5b1e520c2]{Full adder} --- A full adder adds A, B, and carry-in.
\end{itemize}

\section{Used by}\label{appendix-17917d9fdbe3-used-by}

\begin{itemize}
\tightlist
\item
  \hyperref[appendix-b5c46b2d94ff]{Carry-lookahead and parallel-prefix adders} --- uses ripple-carry addition as part of its representation or reasoning.
\end{itemize}

\section{Recall and next depth}\label{appendix-17917d9fdbe3-recall-and-next-depth}

Explain ripple-carry addition without the label, work the example above, then state the assumption whose removal would change the result.

\textbf{Pending:} physical timing, layout, and faster carry networks.

\section{Source and navigation}\label{appendix-17917d9fdbe3-source-and-navigation}

\href{Continuous\%20Notes/Computing\%20Foundations\%20-\%20From\%20Bits\%20to\%20C++\%20and\%20Systems\%20APIs\#Arithmetic\%20becomes\%20a\%20policy}{C002 teaching chapter} · \hyperref[appendix-fd090d982a15]{Part definitions and mechanisms} · \href{../02-Foundations-II/main.pdf\#nameddest=appendix-4ad47bdd269a}{Dependency Map (02-Foundations-II, appendix-4ad47bdd269a)}
